\begin{frame}
    \frametitle{Subespacios vectoriales}

    \begin{definition}
        Sea $V$ un espacio vectorial sobre $k$. Un subconjunto \textbf{no vacío}
        $W\subseteq V$ es un \alert{subespacio vectorial} de $V$ si es
        cerrado para la suma y el producto por escalares:
        \begin{enumerate}
            \item para todo $u,v\in W$, se tiene $u+v\in W$,
            \item para todo $\lambda\in k$ y $v\in W$, se tiene $\lambda v\in W$.
        \end{enumerate}
    \end{definition}

\end{frame}

\begin{frame}
    \frametitle{Subespacios como espacios vectoriales}

    \begin{proposition}
        Todo subespacio $W\subseteq V$ es en sí mismo un espacio vectorial
        con las operaciones heredadas de $V$.
    \end{proposition}

    \begin{proof}
        Tomando $\lambda=0$ y un vector cualquier $v\in W$ en (2), lo cual es posible porque $W$ no es vacío, se obtiene $0\in W$.
        Tomando $\lambda=-1$, se obtiene $-v\in W$ para todo $v\in W$.
        Los axiomas restantes —asociatividad y conmutatividad de la suma,
        distributividades y compatibilidad del producto por escalares—
        valen en $W$ por ser válidos en $V$ para cualesquiera vectores.
    \end{proof}
\end{frame}
